\pdfoutput=1
% ^ MUST stay within the first five lines, and comments count toward that budget. arXiv's TeX
% pipeline decides between latex (DVI) and pdflatex by inspection, and a DVI run cannot
% \includegraphics a PDF figure; without this declaration a submission carrying PDF figures can be
% announced with its figures silently missing.
\documentclass[11pt]{article}

% ---------------------------------------------------------------------------
% Paper III of three: Your Loop Optimizer Is Not Optimizing Your Loop
% Build: pdflatex main && bibtex main && pdflatex main x2
% ---------------------------------------------------------------------------
% Shared preamble for the three papers of this series. Each paper's main.tex declares
% \pdfoutput and \documentclass itself -- \pdfoutput must sit within the first five lines of the
% file arXiv reads, and an \input cannot satisfy that -- then inputs this file for everything else.
%
% Kept as one file rather than triplicated because every setting below was arrived at by fixing a
% specific rendering defect, and three copies would drift apart at the first fix.

\usepackage[utf8]{inputenc}
\usepackage[T1]{fontenc}
\usepackage{amsmath,amssymb}
\usepackage{booktabs}
\usepackage{graphicx}
\usepackage{hyperref}
% Default hyperref draws a coloured rule around every link, so every \S and \ref printed as a red
% box. Invisible in the source and unmissable on paper; found by rasterising a page rather than by
% reading the log, which reports nothing about it.
\hypersetup{hidelinks}
\usepackage{natbib}
\usepackage[margin=1in]{geometry}
\usepackage{xcolor}
% Long \texttt{} identifiers and inline result vectors give TeX nothing to break on, which
% cost eleven overfull boxes running up to 1.7in into the margin. This lets it stretch
% interword space as a last resort instead; content is unchanged.
\setlength{\emergencystretch}{3em}
% Float placement. With many tables queued, LaTeX's defaults (a top float may occupy at most 0.7 of
% the text height, and 0.2 of every float-bearing page must be text) deferred two full-width figures
% to the end of the document --- twenty-eight pages from the section that references them. Relaxing
% the fractions lets a tall figure sit at the top of its own page near its reference instead of
% falling off the back of the float queue.
\renewcommand{\topfraction}{0.92}
\renewcommand{\bottomfraction}{0.75}
\renewcommand{\textfraction}{0.08}
\renewcommand{\floatpagefraction}{0.70}
\setcounter{topnumber}{3}
\setcounter{totalnumber}{4}
\newcommand{\dd}{\ensuremath{\Delta}}


% The PDF's own metadata title, which was empty on every build until 2026-09-03: a reader
% who downloads the file saw `main.pdf` in the window title and reference managers had
% nothing to scrape. Deliberately title only. Setting `pdfauthor` here would put the
% author's name inside the PDF of the ANONYMOUS build too, where the byte scan reads the
% staged sources rather than the compiled metadata, so the blind would break silently.
\hypersetup{pdftitle={Your Loop Optimizer Is Not Optimizing Your Loop}}

\title{Your Loop Optimizer Is Not Optimizing Your Loop\\
\large A silent failure mode in componentized prompt optimization,\\
and the one integer that catches it}

\author{Vladimir Rotariu \\ Independent researcher \\
  \texttt{github.com/loopgepa/whystuck}}
\date{September 2026}

\begin{document}
\maketitle

\begin{abstract}
Reflective optimizers such as GEPA \citep{agrawal2025gepa} can evolve several text components of a
compound LLM system at once. We report a failure mode of that capability which held-out scores
cannot reveal: \textbf{pointing an optimizer at a component is not the same as optimizing it}.
Thirteen of our first fifteen contrast runs measured the effect of evolving a self-refine loop's
verifier against evolving only its generation prompt; the other two pointed the optimizer at a
repair gene instead. Replayed offline, the winning genome's critique gene is
\emph{byte-identical to the seed in $15/15$ cases}, and $13/13$ where it was evolvable at all;
ten of the runs reported held-out scores identical on every seed. The two arms whose difference the study existed to measure were, for the gene that
defines the difference, the same program, and we believed those runs for months. Two ordinary wiring
choices caused it, one a library default: a reflective dataset built without regard to which
component it describes, so the reflector never saw a verdict the critique had returned; and a
component selector that rewrites one component per proposal, so under a matched \emph{total} budget
the two-gene arm searched the shared productive gene half as hard as the one-gene arm, a confound
that flatters the one-gene arm exactly where an effect was being sought. Auditing the fix, the audit,
and a replication, we found the same class of error eight more times, including a ground-truth signal
that was silently constant and an empty completion accepted as a gene, which deletes a step of the
loop and thereby \emph{improves} the training score. We propose a reporting minimum that costs one
integer per component, whether the optimizer's output differs from the seed, and would have
ended this line of work fifteen runs earlier. The failure is not ours alone
\citep{pappu2026multiagent}. Every claim this work withdraws is listed in
Appendix~\ref{sec:corrections}.
\end{abstract}

% The companion-paper note. Every paper in the series carries it so a reader who arrives at one
% knows what the other two contain and, more importantly, which claims are NOT in the paper they
% are holding. Three papers were carved out of one 46-page draft; the carve is documented in
% docs/papers/README.md, and the split's completeness is machine-checked by
% scripts/verify_split_coverage.py against the pre-split source.

\paragraph{This paper is one of three.}
One study produced three separable results, and they are reported separately because they answer
different questions, rest on different experiments, and are useful to different readers. All three
share one harness, one artifact release and one corrections register.
\begin{description}\itemsep0pt
  \item[I. The conditional lever.] Whether optimizing a self-refine loop's verifier beats
    optimizing its generation prompt, and the quantity that decides it: $p_0$, the base rate of
    first drafts already satisfying every requirement.
  \item[II. What a checker cannot buy is the repair.] The loop's own ceiling, with no optimizer
    anywhere in it: how good the intrinsic critique is as a detector, what happens when the
    detector is made perfect, and which fix-side levers actually pay.
  \item[III. Your loop optimizer is not optimizing your loop.] A silent failure mode in
    componentized prompt optimization: pointing an optimizer at a component is not the same as
    optimizing it, the difference does not appear in held-out scores, and one integer per component
    catches it.
\end{description}
Shared artifact, including every cassette, pre-registration and the corrections register:
\texttt{github.com/loopgepa/whystuck}.


\section{Introduction}
\label{sec:intro}

Automatic prompt optimization has matured from single-instruction search
\citep{madaan2023selfrefine} to optimizers that treat an entire LLM program as
the object of optimization. GEPA \citep{agrawal2025gepa}, an ICLR-2026 oral and
the reflective optimizer these runs drive, evolves any set of
text components of a compound system (prompts, rubrics, tool descriptions)
by reflecting on execution traces in natural language and maintaining a Pareto
frontier of candidates. DSPy \citep{khattab2023dspy} similarly compiles and
optimizes multi-module pipelines end-to-end.

That capability creates an experimental hazard with no analogue in single-prompt optimization.
When an optimizer may rewrite one component, ``did it run'' and ``did it change the thing under
study'' are the same question. When it may rewrite several, they come apart: the search can run,
accept candidates, and report a plausible held-out number while the component whose contribution
you are measuring never changes. Nothing in the output distinguishes that case from a genuine
null. A componentized experiment can therefore return a clean, replicated, well-behaved zero from
an experiment that did not happen.

This paper reports such a case in full, on our own artifact, because we did not catch it for
fifteen runs. The study in question asked whether evolving a self-refine loop's
\texttt{ensure\_correct} (verifier) gene alongside its \texttt{generate} gene beats evolving
\texttt{generate} alone. Ten runs reported held-out scores identical on every seed and we read the
consistency as evidence of a tight null. It was evidence of the opposite. Our contributions:

\begin{enumerate}
  \item \textbf{The failure mode, established on a released artifact} (\S\ref{sec:defect}).
    Across the fifteen runs, the winning genome's critique gene is byte-identical to the seed in
    $15/15$, and in $13/13$ of the thirteen that made that gene evolvable at all; and within the
    flagship run, whose five seeds and three evolvable arms give fifteen
    arm-runs, the optimizer changed \texttt{generate} in $5/15$ winners and
    \texttt{ensure\_correct} in $0/15$ ($2/5$ and $0/5$ for the verifier-evolving arm alone). \emph{Two different
    fifteens}, and we name the unit of each because they are numerically equal and an earlier draft merged them. No reported difference measures the
    treatment, and most of them compare a genome against itself, which is why they were so
    stable.
  \item \textbf{Two ordinary causes, one of them a library default} (\S\ref{sec:defect}): a
    component-agnostic reflective dataset, and a one-component-per-proposal selector under a
    matched total budget. Neither is exotic, neither is visible in a held-out score, and the second
    biases the comparison in the direction that manufactures a null.
  \item \textbf{Eight more instances of the same class} (\S\ref{sec:defect},
    Table~\ref{tab:defects}), found by auditing the fix and then auditing the audit. Five
    generalize past our code: a ground-truth signal that is silently \emph{constant} still trains;
    a replay advertised as strict can swallow a cache miss and exit $0$ while reporting a different
    number; an empty LLM completion accepted as a gene deletes that step of the loop, which
    \emph{improves} the training score whenever the deleted step is worse than useless; a sampling
    flag reached one of three sampled surfaces while the metadata recorded it as set; and a pooled
    variance estimator screened every dimension of its regime except which model produced the
    draws.
  \item \textbf{A reporting minimum} (\S\ref{sec:fix}): per component, whether the optimizer's
    output differs from the seed. One integer. Plus a presentation rule for the comparison itself:
    label a difference between two optimized arms with how many seeds \emph{each} arm accepted
    on, which catches the mirror-image defect where the \emph{comparator} is the starved arm.
    A later pre-registered run shows that flagging the zero case is not enough and the contrast
    has to be \emph{stratified} by that count: neither of its arms was starved, yet $74\%$ of its
    positive $\dd$ came from the two seeds where the treatment never applied.
  \item \textbf{A positive control, as a requirement rather than a nicety}
    (\S\ref{sec:posctrl}). An instrument never shown to register a signal cannot be trusted to
    report its absence. On a landscape where the verify gene is the only path to a passing answer,
    the same harness and the real optimizer return a difference of $+1.000$.
  \item \textbf{Measurement noise is a property of the genome, not of the study}
    (\S\ref{sec:sigma}). Measured directly from $135$ re-scorings, $\sigma$ on a thirty-task mean
    ranges over $5.9\times$ across the genomes of one experiment. A single study-wide $\sigma$ is
    an approximation whose direction depends on which genome produced it, and we show a case where
    a borrowed denominator was deciding published claims.
  \item \textbf{The register} (Appendix~\ref{sec:corrections}): the withdrawals this failure mode forced,
    disclosed in one place, including two claims about our own null that turned out to be
    restatements of the defect. The runs it invalidated are kept rather than deleted
    (\S\ref{sec:superseded}), each with the three readings of its \dd\ side by side, since a
    retraction a reader cannot inspect is a claim about a claim.
\end{enumerate}

GEPA is a strong optimizer, and the companion papers of this
series use it as one. The point is that \emph{pointing an optimizer at a component is not the same
as optimizing it}, and that the difference is invisible in exactly the number a reader checks.

\section{Background and Related Work}
\label{sec:related}

\paragraph{Reflective and evolutionary prompt optimization.}
GEPA \citep{agrawal2025gepa} reports beating GRPO by $6\%$ on average and up to $20\%$ while
using up to $35\times$ fewer rollouts, and exceeding MIPROv2 by over $10\%$. GEPA
optimizes \emph{one or more} components and is known to plateau. The sharpest published
instance is independent of us: optimizing a prompt with GEPA across six tasks,
\citet{pappu2026multiagent} report that ``on GPQA Diamond and MATH-500, GEPA returned the
original prompt essentially unchanged.'' That is our failure mode's \emph{signature} in
someone else's pipeline, and reported, which is the practice this paper argues should be routine.
They diagnose no cause, and an unchanged prompt is also what a working search returns when it
finds nothing worth keeping, which is precisely why the signature needs the per-component
integer beside it. OPRO \citep{yang2023opro} and EvoPrompt \citep{guo2024evoprompt} likewise use an
LLM as the optimizer, and OPRO's authors quantify the over-fitting risk directly: their
``training accuracies are often 5\%--20\% higher than our test accuracies''.

\paragraph{More optimizable surface is not reliably better.}
The adjacent question for multi-\emph{agent} systems has been benchmarked:
\citet{bai2026masprompt} find that GEPA lifts single-agent performance by $+4.2$ points
averaged over their nine tasks, with one of the nine flat and one at $-3.3$, while its
multi-agent extension ranges from $-16.0$ to
$+24.0$ points depending on task and workflow topology. Our contribution is orthogonal to that
variance and prior to it: before asking whether componentized optimization helps, one must
establish that it \emph{occurred}. A spread of $40$ points is compatible with a
mixture of real effects and silent no-ops, and nothing in a held-out score separates them.
The same question has been put to the judge prompt directly. Extending textual-gradient
optimization across several evaluation criteria at once, \citet{darshan2026gradients} measure the
gradient's task-focus dropping by $59\%$ when one feedback channel has to serve them all. Their
failure is a loss of \emph{specificity} where ours is a loss of \emph{relevance}, defect~2 of
Table~\ref{tab:defects}, where the reflector read the generator's records while editing the
critique. Both carry the same lesson for a reflective optimizer: what it reads has to be evidence
about the component it is editing, and no score reveals when it is not.

\paragraph{Loop engineering, and why this is a hazard for the practice.}
During June 2026 industry practice converged on a name for designing the machinery
\emph{around} a model rather than its prompt \citep{steinberger2026loops,osmani2026loop,swyx2026loopcraft}.
The stack is commonly decomposed into four loops \citep{langchain2026loopcraft}: an
\textbf{agent} loop calling tools until done, a \textbf{verification} loop scoring output against a
rubric and retrying with feedback, an \textbf{event-driven} loop triggering runs from outside, and a
\textbf{hill-climbing} loop that feeds execution traces to an analysis step which improves the
harness. The last is called ``arguably most important'', and it is the one this paper is about: a
hill-climbing loop improves a harness by attributing trace evidence to components. \emph{If that
attribution is wrong, the optimizer rewrites the wrong thing and every downstream number stays
plausible.} That is what makes the defect below worth a paper rather than a bug report, the
failure is not GEPA-specific but a hazard for any such loop.

\paragraph{Reward gaming, and the second thing to screen for.}
Where the defect is fixed and the optimizer does reach the component, the corner it converges on is
itself worth screening. Structurally this is reward gaming in the sense of
\citet{skalse2022reward}: the training score is a proxy, and a component can be collapsed to a
constant policy the proxy cannot punish. \citet{pappu2026multiagent} report GEPA converging on
\emph{extreme-deference} instructions telling non-expert agents to treat the expert's answer as
correct by default; the companion paper on the conditional lever reports it converging on a verifier that
never accepts anything. Refusing to accept and deferring unconditionally are the same move, which
suggests the corner is a property of reflective search under a weak signal rather than of one task
family. Skalse et al.'s result that unhackable proxies are essentially unavailable is worth stating
here, because it says the defence is not a better training signal but a screen on the optimizer's
output.

\paragraph{Reporting practice.}
Reporting guidance for pre-registered work asks that deviations be disclosed in one systematic place
rather than left for a reader to reconstruct by comparing documents
\citep{willroth2024bestlaid}. Our register (Appendix~\ref{sec:corrections}) follows that guidance, and this
paper's experience adds a specific argument for it: two of our withdrawn claims were not merely
wrong but were \emph{restatements of the defect}, ``the null is tight rather than merely
underpowered'' rested on four intervals of exactly $[0,0]$, which is the sharpest available evidence
that an experiment did not run.

\section{The Study, and What It Was Measuring}
\label{sec:method}

% The harness description common to all three papers of this series. Copied into each rather than
% cited across them, because a reader of any one paper needs it to read that paper's numbers.
%
% Kept as a single input, verbatim, for a reason specific to this study: the claim checker
% (scripts/verify_paper_claims.py) matches several of these sentences as literal text, spaces and
% line breaks included, so three hand-maintained copies would drift and silently stop being
% checked. Edit here, and never reflow a line that contains a number.

\paragraph{The loop.}
The worker is a real generate--critique--refine loop. A \texttt{generate} gene, which the prose also calls the
generation gene, drafts an answer; a \texttt{ensure\_correct} gene, which it calls the
verify, verifier or critique gene interchangeably, drives an LLM
self-critique that returns \texttt{PASS} or a list of fixes and gates a bounded
number of refine rounds; the final answer is scored by a programmatic metric. Two further
genes are optional and named where they are used: a \texttt{plan} module that drafts before
generating, and a repair \texttt{actuator} that turns a critique into the instructions the
refine step is handed.
Critically, the critique is \emph{intrinsic} wherever the loop is left as built: the
programmatic scorer is never consulted inside the loop (that would be oracle contamination), so
the loop cannot game its own metric. Paper~II suspends this deliberately in its oracle-critique
controls, which hand the in-loop gate the true metric and are named as such wherever they are
reported; nothing else in the series does.
Each eval costs \emph{at most} $1 + 2R$ LLM calls for a refine
budget of $R$ (or $2 + 2R$ with the optional plan module), and fewer whenever a critique passes
before the budget is spent: the loop stops on a \texttt{PASS} without issuing the refine call
that would have followed it, so the bound is reached only by an eval that never passes.

\paragraph{Tasks.} A composable IFEval-style \citep{zhou2023ifeval}
multi-constraint generator with nine independent, parameter-carrying verifiers
(e.g.\ \texttt{exact\_word\_count}, \texttt{forbidden\_word},
\texttt{word\_max\_length}, \texttt{num\_bullets}). Joint satisfiability of every
generated task is guaranteed by a witness the generator validates to score
$1.0$ (verified exhaustively: $3{,}189$ easy-tier and $19{,}632$ hard-tier
combinations). The \emph{hard} tier mandates \texttt{exact\_word\_count} (a
genuine LLM skill gap) plus at least one further ``brutal'' constraint, which
structurally caps the seed below saturation.

\paragraph{Metric.} Mean fraction of constraints satisfied on a disjoint
held-out split (\texttt{mc\_score} $\in [0,1]$), scored programmatically (no LLM
judge). We report a relative result under matched budget, not a
leaderboard number.

\paragraph{Reproducibility and budget.}
Every LLM call is recorded to a content-addressed cassette (\texttt{sha256} of the call's name,
positional arguments \emph{and} keyword arguments, which is what gives a per-replicate salt or a
changed temperature its own key rather than the previous draw's), so
runs replay offline at zero cost and partial spend is always persisted. Reproduction is strict:
re-running any study with \texttt{-{}-replay} re-executes the whole pipeline (loaders, GEPA search,
scoring) against the cached calls with no API key, and any cassette miss raises, per call and,
after the correction recorded as defect~7 of Paper~III's defect table, per run as well, since a
consumer that swallows the per-call
error can otherwise let a replay finish with a different number. GEPA is deterministic under
replay, so replay is bit-for-bit, and \texttt{scripts/verify\_reproducibility.py}
registers thirty-eight committed runs. Thirteen are labelled \textbf{superseded} and each
carries its own reason: eleven because the corrected harness changes every reflection prompt and
therefore every cassette key, and two because replaying them re-serves an empty reflection
completion that the blank-gene guard now refuses. The script reports them as skipped rather than
verified.
The remaining twenty-five were recorded under the corrected harness and every per-arm held-out score replays
exactly. The five remaining runs of the forty-three released are the
BigCodeBench ones \citep{zhuo2024bigcodebench}, listed as unreplayable because executing model code re-rolls unseeded
randomness and so shifts the content-addressed prompt key. Those three states partition the
release with nothing left over, and a test asserts they do: a run that supplies a number in
these pages cannot hold no reproducibility status at all. Output tokens are hard-capped via enforced \texttt{max\_tokens}; a pre-flight guard
refuses any run whose conservative estimate exceeds \texttt{-{}-max-usd}, whose default is \$5.
Five runs raised it explicitly, at four distinct values (\$6.18 twice, \$5.89, \$6.52, \$11.76), so
\$5 is a default rather than a study-wide ceiling. Experiments were pre-registered: hypothesis,
frozen configuration, expected seed band and verdict rule committed to version control
\emph{before} spending. Most of them, not all, and the exceptions are named here rather than
left for a reader to reconstruct from commit dates. The first single-constraint run predates the practice and
has no registration. The code-domain runs and the actuator experiment have none either, each
appearing first in the commit that carries its own results. Two further lines, the
feedback-actionability pair and the second-worker replication, do carry registrations, but
committed in the same commit as their results and after the runs' own recorded timestamps, by
seventeen to twenty-five minutes in the first case and by two and a half to three and a quarter
hours in the second. Those registrations may well have been written first; version control cannot
show it, so for them we do not claim it.


\paragraph{Four arms.}
All arms share an identical seed genome; they differ only in what GEPA may
evolve, so a held-out gap isolates componentized optimization from a
``more prompt text'' confound.
\begin{itemize}
  \item \textbf{A}: seed baseline, evolves nothing.
  \item \textbf{B}, prompt-only: evolves the \texttt{generate} gene.
  \item \textbf{C}, command control: evolves the whole genome but scores with
    a genome-blind worker; its landscape is flat by construction (a harness
    soundness check).
  \item \textbf{D}, loop-level: evolves the whole genome
    (\texttt{generate}\,$+$\,\texttt{ensure\_correct}, and \texttt{plan} when
    enabled) with the LLM worker.
\end{itemize}
$\dd = \mathrm{mean}(D) - \mathrm{mean}(B)$ is the quantity of interest.
A and C are expected flat. \textbf{Caveat, established after the fact}
(\S\ref{sec:defect}): as originally implemented neither control could fail, A copied its
own seed score with $\delta$ hardcoded to $0.0$, and C's command loop scored a saturated
$1.0$, so it proposed nothing. The corrected harness re-measures A and gives C a
genome-blind constant draft scored by the real metric, which is flat without being
saturated.

\paragraph{Optimizer.}
We drive the official GEPA \citep{agrawal2025gepa} via a loop-genome adapter, calling
\texttt{gepa.optimize} with \texttt{seed\_candidate}, \texttt{trainset}, \texttt{adapter},
\texttt{reflection\_lm} and \texttt{max\_metric\_calls}, plus a
\texttt{MaxCandidateProposalsStopper} to make the reflection-LM spend a hard cap. GEPA's reflection channel (offline, between
training iterations) sees per-constraint feedback prose: this is the standard
optimizer signal and is identical across B and D; it is distinct from the
in-loop critique gate. Worker \texttt{gpt-4.1-mini} (output cap 320 tokens); reflection
\texttt{gpt-5.6-luna} on the corrected-harness runs this series' conclusions rest on, and
\texttt{gpt-5.6-terra}, which is the CLI default, on the superseded runs that preceded it and on
the later cross-provider critic run (output cap 1500 throughout). The reflection model is therefore \emph{not} constant across the
superseded and corrected sets, so it is one more reason not to read a superseded number against a
corrected one; within each set it is fixed, which is what the contrasts require.

\paragraph{Verdict rule (pre-registered).}
\textbf{SUPPORTED} iff $\dd > \sigma_B + \sigma_D$ and positive;
\textbf{NOT SUPPORTED} iff $\mathrm{mean}(D) \le \mathrm{mean}(B)$;
\textbf{INCONCLUSIVE} if positive-but-within-noise, fewer than two seeds, or a
ceiling/floor flag fires (seed held-out $\ge 0.9$ or $\le 0.1$). This is a
deliberately conservative bar (the sum of two standard deviations), erring
toward under-claiming. The rule has one blind spot: every branch of it reads
$\dd$ and the arm variances, so a genome-against-itself contrast passes through it and returns
\textbf{NOT SUPPORTED}. A verdict rule cannot detect a treatment that never applied.

\section{How a Componentized Treatment Can Silently Not Apply}
\label{sec:defect}

Our first fifteen contrast runs produced arms whose held-out scores were identical on every seed ten times,
and near-zero or moderately negative otherwise. None was a measurement of the hypothesis. Replayed
offline, the winning genome's critique gene is \textbf{byte-identical to the seed in $15/15$} of
those runs. Thirteen of the fifteen made that gene evolvable in arm~\textbf{D} and it is
byte-identical in $13/13$ of them; the other two are the actuator runs, where arm~\textbf{D}
evolved \texttt{generate} and an \texttt{actuator} gene with \texttt{ensure\_correct} held fixed
by design, so identity there is expected rather than diagnostic and what invalidates their
contrast is the selector confound below (defect~3) rather than defect~1.
Within the flagship run, five seeds by three evolvable arms, so fifteen arm-runs,
a second fifteen that we distinguish from the first because the two are equal by coincidence and an
earlier draft ran them together, since GEPA changed
\texttt{generate} in $5/15$ winners and \texttt{ensure\_correct} in $0/15$; for arm~D alone that
is $2/5$ and $0/5$. The two arms whose
difference the study exists to measure were, for the gene that defines the difference, the
same program. No $\dd$ in Table~\ref{tab:results} is therefore a measurement of the
treatment, and most of that column is a genome compared against itself, which is why it was so
stable. Where it is not, the two arms differ in the gene they \emph{share} rather than in the
treated one (\S\ref{sec:superseded}).

Two wiring defects produced that, and neither is visible in a held-out score.
\textbf{The reflective dataset was component-agnostic.} GEPA builds reflection records per
component; ours passed every component the \emph{generator's} inputs and outputs. Asked to
improve a critique gene, the reflector saw drafts and scores but never a verdict the critique
had returned, so it could not observe the behaviour it was editing, and proposed generation prose into a critique slot, which then lost on training score. \textbf{The
component selector rewrote one component per proposal.} GEPA's default \texttt{round\_robin}
cycles, so under a matched proposal cap the two-gene arm spent half its proposals on each gene
and searched the shared productive gene half as hard as the one-gene arm. A matched
\emph{total} budget is not a matched budget per component, and the confound flatters the
one-gene arm exactly where we were looking for an effect.

Auditing the fix, then the audit, then a replication on a second worker,
found the same class of error eight more times.
Table~\ref{tab:defects} is the record. Five are worth the reader's attention because they
generalize past our code, three from the audits and two from the tooling below: a ground-truth signal that is silently \emph{constant} still trains
(defect~6); a replay advertised as strict can swallow a cache miss and exit zero while
reporting a different number (defect~7); and an empty LLM completion accepted as a gene
deletes that step of the loop, which \emph{improves} the training score whenever the step it
deletes is worse than useless (defect~8).

Two arrived later than the rest and are worth attention for a different reason: they were found in
the \emph{instruments}, not in the optimizer, while replicating the ceiling paper's detector result on a
second worker. A \texttt{-{}-temperature} flag was applied where the sampling was assembled for the
worker call and nowhere else, so the critique and verdict calls kept the provider default while the
run's metadata recorded a pinned temperature (defect~10). It surfaced only because two calls issued
on one byte-identical prompt at temperature zero returned different text thirty times out of thirty,
which is not a thing anyone checks unless the answer is supposed to be identical. And the estimator
that supplies this paper's $\sigma$ screened candidate replicates on task set, seed count, levers,
exclusions and gene overrides, every dimension except the \emph{worker} that executed them, so a run on
another model, dropped into the release directory for an unrelated experiment, enrolled as an
eleventh replicate (defect~11). Nothing reached a reader: the claim checker failed on the same
commit and the screen went in before publication, so what follows is a counterfactual, measured
by re-running the checker with the screen removed. It moves $79$ of this paper's published
values, and only $42$ of the $79$ move far enough to breach the tolerance that guards them; the
remaining $37$ shift with nothing objecting. Those figures were $48$ and $29$ until a later audit
added checks for seven $\sigma$-derived values the papers had always printed and nothing had ever
read, then $55$ and $30$ when a further one added an eighth, then $56$, $57$, $60$, $68$ and now
$79$ as we scoped Paper~I's superlative, drew its figures, unified the noise floor's printed
precision, and swept the replicate count across all three papers, each new check turning out to
read
$\sigma$ too: the contamination was moving them all
along, and the count of what it moves is bounded by
the count of what is watched. The pair has now moved seven times for that reason and not once
because the contamination changed, which is the cleanest available demonstration of the sentence
it sits in. That is the part worth keeping, and it is the same
threshold problem as defect~6 wearing different clothes: a contamination does not have to be
large to be real, only larger than whatever is watching. Both are the paper's own
subject arriving in its tooling: a setting reported as applied that a surface never received, and a
guard reporting agreement about a dimension it never looked at. Neither needed an optimizer to
happen, and each was caught by a machine rather than by reading.

\begin{table}[h]
\centering
\small
\begin{tabular}{p{0.035\textwidth}p{0.35\textwidth}p{0.49\textwidth}}
\toprule
\# & Defect & Why held-out scores could not reveal it \\
\midrule
1 & Critique gene never changed: $0/15$ winning genomes & The arms differ only in which genes
  are \emph{evolvable}, not in which changed; both report a score \\
2 & Reflective dataset component-agnostic (root cause of 1) & The reflector still returns
  fluent, plausible prose; only its \emph{relevance} is lost \\
3 & \texttt{round\_robin} halved the two-gene arm's search per gene & Total budget matched, so
  the comparison looks fair by construction \\
4 & Exact-resample interval dropped the largest atom, which at $n{=}3$ carries
  $1/27 = 3.7\% > \alpha/2$ & Error direction flatters negative effects: brackets zero
  $\to$ appears to exclude it (Table~\ref{tab:ciexcl}) \\
5 & Seeds re-partition one fixed task pool rather than resampling it & Across-seed spread
  looks like replication variance and reads as precision \\
6 & The fix's own ground-truth signal was a \emph{constant} & A constant label is a valid
  label; training proceeds and reports normally \\
7 & A ``strict'' replay swallowed a cassette miss, exiting $0$ with a divergent number &
  Exit code $0$ plus a plausible report is indistinguishable from success \\
8 & The harness accepted an empty reflection completion as the critique gene & A loop with \emph{no}
  critique beats one with a near-random critique on the training subsample, so the degenerate
  candidate wins on merit \\
9 & The \emph{comparator} arm accepted nothing on $3/3$ seeds, so $\dd(D-B)$ was
  $\dd(D-\mathrm{seed})$ & Both arms report a genome and a score; a genome equal to the seed is
  indistinguishable from one the search chose to keep \\
10 & \texttt{-{}-temperature} reached only the \emph{worker} call; critique and verdict kept
  sampling at the provider default & The metadata records the pinned temperature, so the artifact
  asserts a setting two of its three sampled surfaces never got \\
11 & The pooled-$\sigma$ estimator screened task set, seeds, levers, exclusions and gene
  overrides, but not the \emph{worker}, so a run enrolled from another model & A pooled sd moves
  smoothly when a foreign draw joins: every ratio on it shifts a little and none looks wrong \\
\bottomrule
\end{tabular}
\caption{Eleven ways a componentized-optimization result can be confidently wrong while every
number it prints looks reasonable. Defects 1--3 invalidated the original headline; 4--5 change
what the intervals may be read to mean; 6--7 were caught before spending on the re-run and 8
only after it had already corrupted one seed of a paid run; 9 was caught after spending, in a later
sweep, and is defect~1's mirror image, the same one-sided contrast with the starvation on the
control side instead of the treatment side. 10--11 arrived while replicating on a second worker
and are this paper's own thesis turned on its instruments: a switch reported as set that a
surface never received, and a guard reporting agreement about a dimension it never compared.
Neither is specific to an optimizer. 11 was caught before publication; removing its screen again moves $79$ of this paper's values, only $42$ of them past the tolerance that guards them.}
\label{tab:defects}
\end{table}

\begin{table}[h]
\centering
\small
\setlength{\tabcolsep}{5pt}
\begin{tabular}{lcp{0.20\textwidth}p{0.24\textwidth}}
\toprule
Result & per-seed & as published & corrected \\
\midrule
Actuator, powered ($n{=}3$) & $-0.15, -0.15, 0.00$ & $[-0.150,-0.050]$ ``excludes zero''
  & $\mathbf{[-0.150,+0.000]}$ \emph{includes zero} \\
Actionability, hi$-$lo ($n{=}3$) & $+0.090, +0.025, +0.075$ & $[+0.025,+0.085]$
  & $[+0.025,+0.090]$ (still excludes zero) \\
Null, run~3 ($n{=}5$) & $0,+0.055,0,0,-0.055$ & $[-0.033,+0.033]$ & unchanged \\
\bottomrule
\end{tabular}
\caption{The upper bound was too low by one atom. Only $n{=}3$ intervals are affected, and
the direction of the error flatters negative effects: it turns an interval that brackets
zero into one that appears to exclude it. Run~3 is the flagship hard-multi run, listed with
the rest of the numbering in Table~\ref{tab:results}.}
\label{tab:ciexcl}
\end{table}

\section{The Fix, and the Reporting Minimum We Propose}
\label{sec:fix}

The corrected harness (i) builds a \emph{component-specific} reflective dataset, so the verify
gene's records carry the draft it judged, the verdict it returned, and whether that verdict
was a false pass or a false alarm; (ii) passes an explicit component selector (\texttt{module\_selector}), defaulting to
rewriting all genes per proposal so a multi-gene arm is not starved; (iii) replaces arm~C
with a genome-blind constant draft scored by the real metric, which is flat without being
saturated, and re-measures arm~A instead of copying it; and (iv) records, for every arm, the
per-gene changed-from-seed count, the train aggregate, and the final gene text. That count is
an identity test on the winner, not a tally of accepted proposals: GEPA proposed critique
rewrites in the flagship run and they lost on training score, which the released summary records
and which an empty count cannot distinguish from a search that never proposed anything. The audit of
that harness (defects~6--7) added three more: (v) the verify gene's truth signal now receives
the real task, and a trace that cannot be matched to an evaluated task raises rather than being
scored against a stand-in; (vi) a replay that swallows a cassette miss now fails at context
exit instead of exiting $0$ with a divergent report; and (vii) every run releases the
\emph{per-task} held-out score vector behind each mean, which is what makes a paired
task-level test possible for readers who (correctly) do not want to treat our seeds as
replications. Designing a later split-size sweep added an eighth, of the
same shape: (viii) a configuration whose metric-call budget cannot fund a single accept cycle
at the requested train-split size is now \emph{refused} at the dry-run step rather than
priced, because such a run bills in full and reports a flat $\dd$ with
\texttt{genes\_changed=[]}, a null-shaped result from an arm that never searched. That is
this paper's central failure mode arriving as a \emph{budget} arithmetic slip rather than a
wiring one, which is the argument for making the refusal a property of the harness instead of
a habit of the operator. Running that sweep then produced a ninth (defect~9), which is the
same failure mode again from the one direction we were not watching: at
$n_{\mathrm{train}}{=}4$ the \emph{comparator} arm accepted no proposal on any of three seeds,
so its released genome is byte-identical to the seed and $\dd(D-B)$ on that row is
$\dd(D-\mathrm{seed})$. The guard added as~(viii) does not catch this and cannot: it enforces
that an arm \emph{can} fund an accept cycle, never that the search accepts one. Both are
one-sided contrasts wearing a two-sided label, and the difference is only which arm was
starved: in the retraction it was the treatment, here the control. The reporting fix is the
same number~(iv) already demands, so what we add is not a guard but a \emph{presentation} rule,
now enforced by the analysis script: \textbf{a $\dd$ between two optimized arms must be
labelled with how many seeds each arm accepted on}, and flagged where either is zero.
Sixteen released runs trip that flag. Exactly one is \emph{one-sided} in the new sense, with a
starved comparator facing an arm that did search: the $n_{\mathrm{train}}{=}4$ row. In the other
fifteen \emph{both} arms count zero, but by an \emph{absent} field rather than a measured one:
those runs predate item~(iv), so no arm record carries \texttt{genes\_changed} at all and the
census reads the missing key as a zero: this paper's own failure mode, inside the diagnostic we
propose as the fix. Their released scores say otherwise. An arm whose search accepts nothing
returns the seed genome and therefore reproduces its own seed score, so any seed on which
$\mathrm{opt} \ne \mathrm{seed}$ carries an accepted proposal: \textbf{eleven} of the fifteen
have one, arm~B in eleven runs and arm~D in ten, and only four are flat on both arms on every
seed. Figure~\ref{fig:noop} puts the three readings of those fifteen runs side by side, which is
the shortest way to see that a missing field and a measured zero are different findings.
Those fifteen are reported as uninformative rather than as evidence
of no effect, but for the narrower reason that their acceptance record is missing, not that it
is zero.

\paragraph{The rule, applied prospectively, and found too coarse.}
Everything above is retrospective: a rule derived from runs that had already gone wrong. A later
pre-registered run applied it in advance, and we found the rule as stated is not
enough. The run is a cross-provider critic study (Paper~I, D3) over \textbf{eight} seeds; its
acceptance census reads arm~B on \textbf{five} seeds and arm~D on \textbf{three}, cross-checked
against the seed-score identity. Neither count is zero, so \emph{the flag does not fire}. Its
$\dd(D-B) = +0.014$ would be read as a small positive.
Stratifying the same contrast by which arms accepted says something else:

\begin{center}\small
\begin{tabular}{lccr}
\toprule
stratum & seeds & mean $\dd(D-B)$ & share of the total \\
\midrule
arm~D accepted, the treatment applied & 3 & $+0.0100$ & $26\%$ \\
\textbf{only arm~B accepted, D returned the seed} & \textbf{2} & $\mathbf{+0.0425}$ & $\mathbf{74\%}$ \\
neither accepted & 3 & $+0.0000$ & $0\%$ \\
\bottomrule
\end{tabular}
\end{center}

\noindent
\textbf{Three quarters of the positive contrast comes from two seeds on which the treatment did not
apply and the comparator damaged itself}, losing $0.065$ and $0.020$ against its own seed genome.
On those seeds $\dd(D-B)$ is $\dd(\mathrm{seed}-B)$: the same one-sided contrast as
defect~9, at the level of the individual seed rather than the run, and therefore invisible to a
flag that asks only whether an arm's count is zero. Restricted to the seeds where arm~D's winner
moved, $\dd = +0.0100$ on $n = 3$. That restriction conditions on an outcome of the treatment
rather than on the assignment, so $+0.0100$ is not a treatment-effect estimate and we do not
report it as one: what the stratification establishes is where the aggregate comes from, which
is a reason to distrust the aggregate rather than a better estimate to put in its place.

So the presentation rule needs strengthening, and this is the form we now propose: \textbf{report
the per-arm changed-from-seed count per seed, and stratify the contrast by it}, not merely flag the
case where a count is zero. A binary flag catches the run that never searched; it does not catch
the run that searched on some seeds and not others, which is the ordinary case once a budget is
tight enough to matter. The cost is the same one integer item~(iv) already demands, printed per
seed instead of per run.

Item~(iv) is the part we would urge on others: \textbf{a
componentized-optimization result should report, per component and per seed, whether the
optimizer's output differs from the seed.} Had we printed that one number, this study
would have been a different study fifteen runs earlier, and, as the stratification above
shows, per component alone would have been enough to catch the retraction but not the
cross-provider run. Because
the corrected wiring changes every reflection prompt, it changes every cassette key, so the
released cassettes cannot replay under it; they are retained and labelled superseded. The
first runs recorded under the corrected harness are two hand-written-gene conditions reported in
the companion paper on the conditional lever; they replay bit-for-bit, so \texttt{verify\_reproducibility.py}
verifies twenty-five runs rather than reporting that nothing has been proven, those two
conditions, the four-arm re-run, the two new levels of that paper's $p_0$
axis and the two calibration passes behind
them, the bought rows of the split and budget sweeps, the cross-provider critic run, the
eight runs of the second-worker line and the diagonal splice control. The cross-provider run is the only one of them whose critique
was served by a second vendor and therefore the only one that would notice a replay quietly
falling back to the worker's provider. Eleven of those twenty-five were replaying already and
were not registered, which is its own instance of this paper's subject: an unstated
status is read as the stated guarantee covering it.
The \emph{discarded first attempt} at the four-arm re-run is not among them,
for a reason worth recording: the blank-gene guard that attempt motivated (defect~8) makes
its own cassette unreplayable, since a replay re-serves the empty completion and the guard
raises on it. Its numbers stand as measured and are no longer verifiable under the current
code, which is the honest definition of superseded.

\paragraph{Why the diagnostic has to be per component, and printed.}
Two weaker versions of this rule do not work, and both were tried here. Reading
\texttt{genes\_changed} as a whole is not enough: in $5/15$ of our winning genomes it was
non-empty, because \texttt{generate} had changed, so a reader checking ``did the optimizer do
anything'' saw a yes on a run where the treatment gene had not moved. And inferring the answer
from the score is what fails, the tighter the reported null, the stronger the
apparent result and the more likely it is an identity. The rule therefore has to be indexed by
component and printed in the artifact, not derived on demand, because the person who needs it
most is the reader who has no reason to suspect anything is wrong.

\section{A Positive Control: the Instrument Does Respond}
\label{sec:posctrl}

A null is worth reading only if the instrument that produced it can register a signal, and
after three silent failures we do not ask the reader to take that on trust. We built a
landscape in which evolving the verify gene is the \emph{only} path to a passing answer: the
generator always drafts \texttt{x\,x}, which fails the task's single constraint however its
prompt is worded; the critique catches the flaw only when the verify gene carries a sentinel;
and a caught flaw triggers a refine that passes. Arm~B therefore cannot escape the failing
draft no matter how well it optimizes, and arm~D can. Running the real four-arm study through
real GEPA on this landscape (no stubbed optimizer, no network, no spend) gives
A~$0.000$, B~$0.000$, C~$0.000$, D~$\mathbf{1.000}$, i.e.\ $\dd(D-B) = +1.000$. The harness
detects a loop-level advantage when one exists, so the near-zero $\dd$ reported on real tasks
is a property of those tasks and that optimizer, not of the measurement path. The control is a
saturating one, $0.000$ against $1.000$, so it establishes that the path is not \emph{dead}; it
says nothing about sensitivity at the $0.01$--$0.04$ scale the real contrasts live at, against a
$\sigma(\dd)$ near $0.03$. A canary is not a calibration
(\texttt{scripts/positive\_control.py}). It also scores its four tasks as both train and
held-out, which the shared method's disjoint-split rule does not permit and which a canary does
not need: the question it answers is whether the path can register a signal at all.

We state this as a requirement rather than a courtesy. The control is cheap: it needs no
network and no spend, and it is the only artifact in this study that would have failed
loudly under the original wiring. Had we built it first, it would have caught defect~1 before
any money was spent, because on that landscape a control that cannot fail shows itself: the
command arm scores $1.000$ where a genome-blind worker has to score $0.000$. It would not have
caught defects~2 or~3, and we say so rather than claiming the control is worth more than it is.
Its reflection function returns the same proposal whatever records it is handed, so a
component-agnostic reflective dataset is invisible to it; and it runs to this day with
\texttt{module\_selector} at GEPA's default, so the selector confound is present in the control
itself and the control still returns $\dd(D-B) = +1.000$. An instrument that has never been shown to register a signal
cannot be trusted to report its absence, and a componentized optimizer is exactly the kind of
instrument whose failure is silent.

\section{Superseded Runs: What the Defect Invalidated}
\label{sec:superseded}

The original ten runs below are retained as the record of what the defect of
\S\ref{sec:defect} invalidated, and are \textbf{not} evidence about loop-level
optimization: the loop-level gene never changed, so nothing in the $\dd$ column measures the
treatment. Most of it compares a genome against itself; the rows where it does not are the
selector confound read through the shared gene, which the discussion below separates. Their \textbf{Seed} column is unaffected: it
scores a fixed genome and involves no optimizer, and it is the live readout the companion
paper on the loop's ceiling draws on for its oracle, fix-budget and actionability results. The
cassettes are released and labelled superseded: the corrected wiring changes every reflection
prompt, hence every cassette key, so these runs cannot replay under current code.

\begin{table}[t]
\centering
\small
\setlength{\tabcolsep}{4.5pt}
\begin{tabular}{clccccl}
\toprule
Run & Regime & Seeds & Seed & B (prompt-only) & D (loop-level) & \dd \\
\midrule
1 & single-constraint        & 2 & 0.47 & 0.59 & 0.50 & $-0.09$ \\
2 & easy multi ($k{=}3{-}4$) & 5 & \textbf{0.97} & 0.97 & 0.97 & $0.00$ \\
3 & hard multi ($k{=}4{-}5$) & 5 & 0.765 & 0.770 & 0.770 & $\mathbf{0.00}$ \\
4 & \quad $+$ feedback tool  & 3 & 0.837 & 0.798 & 0.798 & $\mathbf{0.00}$ \\
5 & \quad $+$ plan module    & 3 & 0.757 & 0.757 & 0.757 & $\mathbf{0.00}$ \\
6 & \quad $+$ cross-lineage critic & 3 & 0.825 & 0.817 & 0.825 & $+0.01$ (INC.) \\
7 & \quad $+$ oracle critique ($R{=}1$) & 3 & \textbf{0.748} & 0.74 & 0.74 & $\mathbf{0.00}$ \\
8 & \quad $+$ oracle critique ($R{=}3$) & 3 & \textbf{0.813} & 0.80 & 0.80 & $\mathbf{0.00}$ \\
9 & \quad $+$ oracle ($R{=}3$), feedback \emph{lo} & 3 & \textbf{0.810} & 0.80 & 0.80 & $\mathbf{0.00}$ \\
10 & \quad $+$ oracle ($R{=}3$), feedback \emph{hi} & 3 & \textbf{0.873} & 0.86 & 0.86 & $\mathbf{0.00}$ \\
\bottomrule
\end{tabular}
\caption{All ten of the original runs, retained as the record of what the defect
invalidated. \textbf{The \dd\ column is not a measurement} and no claim in this series
rests on it: the loop-level gene never changed, so nothing here measures the treatment
(\S\ref{sec:defect}), and where the column is non-zero the arms differ in the gene they share
rather than in the treated one; the companion paper on the conditional lever measures this contrast on a corrected harness.
Rows 4--10 elide a repeated ``hard multi'' prefix: each is run 3's regime plus the listed lever. The \textbf{Seed} column is live throughout, and is the intended readout for runs 6--10.
Runs 1--2 are uninformative by design (the verify gene is inert on single-constraint
tasks; the seed saturates at $0.97$ on easy multi). Runs 8--10 share the oracle $R{=}3$
regime and vary only repair-feedback actionability (mid $=$ run 8): terse verdict (lo,
$0.810$) $\approx$ which-constraint (mid, $0.813$) $\ll$ concrete measured edits (hi,
$0.873$).}
\label{tab:results}
\end{table}

\begin{figure}[t]
\centering
\includegraphics[width=\textwidth]{figures/fig-noop.pdf}
\caption{The pre-correction contrast runs, read three ways, and only the middle reading says
anything about acceptance. \emph{Left:} the per-component integer this paper proposes is
\emph{absent} from all \textbf{fifteen} files, since the field postdates them, so a census that
reads a missing key as a zero reports that neither arm ever accepted a proposal. \emph{Middle:} the
released scores contradict that in \textbf{eleven} of the fifteen. An arm that accepts nothing
returns the seed genome and reproduces its own seed score, so a moved score is an accepted
proposal; the converse does not hold, which makes eleven a lower bound rather than a count.
\emph{Right:} the \dd\ column these runs published, against $\pm\sigma(\dd)$. \textbf{Ten} scored
identically on every seed, which is what one program measured against itself looks like in data.
The filled marker at zero is \texttt{multihard}, whose per-seed differences cancel to
$-2.22 \times 10^{-17}$. Rows carry the released artifact names under \texttt{docs/}, and the
\dd\ column is not a measurement (\S\ref{sec:defect}).}
\label{fig:noop}
\end{figure}

\paragraph{What the \dd\ column does not say.} We report it once, above, and do not
interpret it further. In every run in Table~\ref{tab:results} the loop-level arm's extra gene
was never changed (\S\ref{sec:defect}), so no run here measures the treatment, and the four
exactly-$[0,0]$ intervals are evidence that it was not applied rather than evidence that it
did nothing. We previously called those intervals ``tight, not merely underpowered'' and that
reading is withdrawn.

A second reading needs withdrawing with it. We
described the zeros as \emph{identities}, arms B and D issuing identical calls, so the
cassette served identical completions. That is too strong, and its falsity is visible in this
table: run~1 reports $-0.09$ and run~6 $+0.01$, and run~3's per-seed vector is
$[0,+0.055,0,0,-0.055]$, so on several seeds the two arms demonstrably did \emph{not} issue the
same calls. An unchanged verify gene does not make the arms the same program, because they
still differ in the gene they \emph{share}: under the default one-component-per-proposal
selector, arm~D searched \texttt{generate} about half as hard as arm~B (defect~3). So the
non-zeros in this column are not residual treatment effects; they are the selector confound
measured through the shared gene, plus sampling noise, which this study can now put a number on.
Of the
fifteen contrast runs recorded under the pre-correction harness, ten have arms whose held-out scores are
identical on \emph{every} seed and four have $|\dd| \ge 0.005$, the largest being the powered actuator run at
$-0.100$. The identities are not noise measurements: B and D produced byte-identical genomes
there, so they issued identical calls and the cassette served identical completions.

That is ten of fifteen, and the remaining one is worth naming rather than rounding into either
group, because it is a sharper version of this paper's own point. \texttt{multihard} has per-seed
differences $[0, +0.055, 0, 0, -0.055]$: the arms demonstrably \emph{did} differ, on two seeds, and
the differences cancel. Its \dd\ is therefore zero without being an identity, a zero that hides
not one non-experiment but two offsetting effects. We previously reported this run among the exact
zeros, which was true of the number and false of the reason. The four
non-zeros are a different matter again, and must be read against the noise floor
\S\ref{sec:sigma} measures. Scaled by that $\sigma(\dd)$, two of the
four sit under $1\sigma$, the worker swap at $-0.014$ and the cross-\emph{lineage} run at $+0.008$,
and two do not: the powered actuator run at $\approx\!3.44\sigma$ and run~1, whose
single-constraint tasks leave the verify gene inert, at $-0.094$, $\approx\!3.23\sigma$. The caveat
is that $\sigma$ was estimated on
the hard-multi tasks while the largest of these is a code run, so its transfer is an assumption and not a
measurement. This is the same conflation the companion paper on the conditional lever finds in the corrected re-run, one
harness generation earlier: \dd\ is a whole-system contrast, so it moves when the shared gene
moves, whether or not the treatment ever applied.

\paragraph{Two factual slips belong with the withdrawals.} Both were stated about this table and
both were wrong in the direction of understating the damage:
the reproducibility script registers thirteen superseded runs, not ten,
and the \$5 spend ceiling is a CLI default that five runs overrode
(up to \$11.76) rather than a study-wide cap. Neither changes a result; both are the kind of quiet
arithmetic drift a reader finds and an author does not, which is why they are in the register
(Appendix~\ref{sec:corrections}) rather than silently fixed.

What survives in these runs is the \textbf{arm-A seed baseline}, which is a measurement of
the loop rather than of the optimizer and is untouched by the defect. Read down the ``Seed''
column: it is the readout for runs 6--10. Run
6's critic is worth one line of it: routing the critique to \texttt{gpt-5.6-luna} over a
\texttt{gpt-4.1-mini} worker lifted the baseline $0.765 \to 0.825$, directionally consistent
with the reliable-critic hypothesis, but \texttt{gpt-5.6} and \texttt{gpt-4.1} share a vendor
and training lineage, so this is cross-\emph{lineage}, not the cross-\emph{provider} test
that would rule out self-preference bias.

\section{\texorpdfstring{$\sigma$}{Sigma} Is a Property of the Genome, Not of the Study}
\label{sec:sigma}

The defect above is about whether an experiment ran. This section is about what its numbers may
be divided by, which is a second way a componentized study goes quietly wrong: a study that
reports every effect in units of ``its'' measurement noise is asserting that one such quantity
exists.

\paragraph{The estimate that fell out of the artifact.}
No interval in this study contains within-condition sampling
variance, because the worker runs at the provider's default temperature with one sample per
task, so nothing was measured twice on purpose. It was nonetheless measured by accident. Ten
released runs, a four-arm re-run and its discarded first attempt (defect~8; its seed column
involves no optimizer, so excluding it from the \emph{verified} set does not exclude it here),
each row of a split-size sweep and each bought row of a
budget sweep, evaluate the
\emph{same} unoptimized seed genome on the \emph{same} thirty held-out tasks before optimizing
anything, so they are ten independent draws of one fixed number. Their means span $0.7400$ to
$0.8183$: a range of $0.078$ on a quantity that is constant by construction, which is the
plainest statement of the problem this section exists to quantify.
From the paired differences, $\sigma$ for one task under one draw is $0.1125$; a $\dd$ compares
two independent draws, so $\sigma(\dd)$ on one task is $0.1591$ and on a thirty-task mean
\[
\sigma(\dd) \;=\; 0.1591/\sqrt{30} \;=\; 0.0291 .
\]
The estimator is worth stating as a recipe, because ``from the paired differences'' underdetermines
it and an independent auditor tried four conventions without hitting the figure: pair the ten draws
all $45$ ways, take the \emph{sample} standard deviation of the thirty per-task differences inside
each pair, and average those forty-five numbers. That average is the $0.1591$; the $0.1125$ is it
over $\sqrt{2}$, so the paired quantity is the one measured and the single-draw one is derived.
\texttt{scripts/analyze\_measurement\_noise.py} prints all three.
Two caveats travel with it: the mean signed difference
between replicates, $-0.014$, is confounded with any provider-side drift between recording
dates, and five replicates of three seeds cannot separate the two. It establishes an order of
magnitude, and that order of magnitude is the same as the effects being reported.

\paragraph{Then we measured it directly, and it was not one number.}
We bought the direct version: every released genome of the $p_0$ sweep re-scored $K=5$ times
on its own held-out tasks, three $p_0$ levels ($L_1$ to $L_3$, the conditional-lever paper's
constraint-exclusion ladder from $p_0 = 0.067$ to $0.667$) $\times$ \{prompt-only, loop-level,
unoptimized\}
$\times$ three seeds $\times$ five replicates, $135$ runs. No re-optimization is involved, so
this is the existing scoring path on a fixed genome, and all $135$ replay from their cassettes
offline. Two properties make it a measurement rather than a cache: repeated scorings of one
genome are \emph{identical} calls, so a per-replicate salt enters the content-addressed key
without reaching the provider (without it, record-mode dedup would hand $K$ replicates one
cached draw and report zero variance); and no temperature is pinned, because the point is the
noise in the numbers the paper already reports, which were drawn at the provider's default.

\begin{table}[h]
\centering
\small
\begin{tabular}{lccc}
\toprule
level & prompt-only (B) & loop-level (D) & unoptimized (seed) \\
\midrule
$L_1$ & 0.0326 & 0.0230 & 0.0161 \\
$L_2$ & \textbf{0.0376} & 0.0170 & 0.0157 \\
$L_3$ & 0.0159 & \textbf{0.0064} & 0.0111 \\
\bottomrule
\end{tabular}
\caption{Sd of one draw of a thirty-task mean, by genome. The range is $0.0064$ to $0.0376$,
a factor of $5.9$. Two mechanisms are visible: a genome nearer the metric ceiling has less room
to vary ($L_3$'s loop-level arm scores ${\approx}0.98$), and a critique that refines
\emph{every} draft issues more sampled calls per task than one that accepts some. So ``the
noise'' is not a constant of the study, and a single $\sigma$ applied study-wide is an
approximation whose direction depends on which genome it came from.}
\label{tab:sigmagenome}
\end{table}

That splits the section's own $\sigma$ into two quantities that were being conflated, and it is
worth being explicit about which belongs where. A comparison \emph{between two runs recorded on
different dates}, the seed-baseline curve of the companion ceiling paper being one, must carry
provider-side drift, and the cross-run estimate is the one that contains it, so it is the right
denominator \emph{of the two} for that curve. It is not the last word on it: the ceiling paper has
since re-scored each of those rows five times and superseded its ratios with each row's own measured
$\sigma$, which is better than choosing correctly between two borrowed ones. A comparison \emph{between two arms
inside one run} carries no drift, and there the directly measured value applies: $\sigma(\dd)$
is $0.047$ at $L_1$, $0.046$ at $L_2$ and $0.022$ at $L_3$. Two of the three exceed the
$0.0291$ the cross-run estimate gave, because optimized arms vary more than the seed genome that
estimate was built from; $L_3$ does not, which is the ceiling effect
Table~\ref{tab:sigmagenome} shows directly: at $p_0 = 0.667$ both arms sit near $0.97$ and
have little room to vary.
On a \emph{single} draw the three $p_0$-sweep contrasts of the conditional-lever paper,
taken in the same $L_1$, $L_2$, $L_3$ order as the denominators above, are therefore $0.89\sigma$,
$1.31\sigma$ and $0.08\sigma$: weaker against the within-run denominator for two of them,
and marginally stronger for the third, which is a null either way.

One control, without which the table is void: the unoptimized genome's thirty-task mean moves
across replicates with sd $0.0161$, comfortably non-zero, so the salt is producing independent
draws rather than five copies of one cached call. That number is also a check on the offline
estimate, which put the same quantity at $0.0205$: the cross-run figure is $1.3\times$ larger,
which is the size of the drift contamination its own caveat predicted.

\paragraph{The consequence for reporting.}
A borrowed $\sigma$ can decide a claim, and this study has a worked case. Six ratios in the
companion ceiling paper divided gaps by the cross-run $0.0291$, estimated on a \emph{seed}
genome at $R{=}1$ in the $p_0$ family; three of the six had a margin against the $2\sigma$ line
smaller than the $5.9\times$ spread above. Measured on their own rows, every one of those rows
turned out \emph{quieter} than the borrowed denominator, between $0.38$ and $0.82$ of it, so the
borrowing had been conservative and one claim had been withdrawn that should not have been. The
direction was not knowable in advance and we had guessed it backwards. So the reporting rule is
narrow and cheap: state which genome a $\sigma$ was measured on, and do not divide a gap by a
$\sigma$ from a different one without saying so.

\section{Limitations}
\label{sec:threats}

\textbf{This is a case study, not a survey.} One optimizer, one adapter, one task family. We
cannot say how often the defect occurs in the wild; we can say that two of its causes are
ordinary, that one is a library default, that it survived our own review for fifteen runs, and
that an independent group reports a compatible observation \citep{pappu2026multiagent}. The
prescription (\S\ref{sec:fix}) is justified by its cost rather than by an estimated base rate:
one integer per component is cheap enough that it does not need one.

\textbf{The diagnostic is necessary, not sufficient.} A per-component change count catches a
treatment that never applied. It does not catch a treatment that applied and did something
degenerate: the conditional-lever paper's non-accepting verifier changes the gene on every seed. That
needs its own screen, and the cheap one there is a \texttt{PASS} rate of zero. Worse, it reports
the \emph{reassuring} answer on defect~8: an empty completion accepted as a gene \emph{does}
differ from the seed, so the count says ``changed'' on the run where the change deleted
a step of the loop and won on the training subsample by doing so. The blank-gene guard covers that
case and the count does not, which is why we present the two together rather than the integer
alone. Nor does the count
distinguish a semantically meaningful rewrite from a cosmetic one; it is a floor on evidence, not
a measure of treatment strength.

\textbf{The eleven defects are not independent findings.} They were found by auditing one codebase
in sequence, so their number reflects how long we looked rather than a density. What we claim is
the \emph{class}, and that its members share a signature: the failure leaves the reported number
plausible.

\textbf{The $\sigma$ result is one family of genomes, and the cells are not single genomes.}
The $5.9\times$ spread is measured across nine arm-by-level cells of one experiment, each of
which pools its three seeds' genomes before the replicate means are taken, and the three levels
differ in which constraints are excluded. So the spread is across cells that differ in genome,
seed and task regime at once, not across nine genomes holding the rest fixed. The companion
ceiling paper measures $2.2\times$ across six of its own. Both refute a single study-wide constant; neither establishes the size of the
spread as a general quantity.

\textbf{We cannot re-run the superseded cassettes.} The corrected wiring changes every reflection
prompt and therefore every cassette key, so the fifteen invalidated runs are documented rather
than re-analysed. Their numbers stand as measured and are not verifiable under current code, which
is what we mean by superseded.

\section{Conclusion}
\label{sec:conclusion}

A componentized optimizer can be pointed at a component, run, accept candidates, report a
plausible held-out number, and leave the component exactly as it found it. Our first fifteen
contrast runs did
exactly that: the critique gene is byte-identical to the seed in $15/15$ winning genomes, ten runs
reported held-out scores identical on every seed, and we read the consistency as a tight null for months. Two
ordinary wiring choices sufficed, a reflective dataset built without regard to which component
it describes, and a selector that rewrites one component per proposal under a matched total budget,
which is a library default and biases the comparison toward the one-gene arm. Auditing the fix, the audit and a
replication found eight more errors of the same class, whose common signature is that the number they produce
looks reasonable.

The generalizable output is one integer: \textbf{per component, whether the optimizer's output
differs from the seed}. It would have ended this line of work fifteen runs earlier. It has to be
per component, because a whole-genome ``did anything change'' answered yes on $5/15$ of our runs;
and it has to be printed in the artifact, because the reader who most needs it is the one with no
reason to suspect a problem. Three companion rules cost as little. Run a positive control before
spending, since an instrument never shown to register a signal cannot be trusted to report its
absence, and ours would have caught defects~1--3 for free. Label a difference between two optimized
arms with how many seeds \emph{each} arm accepted on, which catches the mirror-image case where the
comparator is the starved arm. And state which genome a $\sigma$ was measured on, because $\sigma$
ranges over $5.9\times$ within one experiment and a borrowed denominator decided, and in our
case wrongly withdrew, a published claim.

There is a broader reading for anyone building a hill-climbing loop over their own harness. Such a
loop improves a system by attributing trace evidence to components \citep{langchain2026loopcraft};
if the attribution is wrong, it optimizes the wrong thing and every downstream number stays
plausible. The defect is therefore not a property of GEPA, which is a strong optimizer, but of the
seam between an optimizer and the system it is pointed at, and that seam is exactly what
componentized optimization adds.

We release the corrected harness, every cassette that has a released result behind it, and
every pre-registration, including the superseded ones, and the audit that produced this paper, on the view that a study which invalidates its own
headline is more useful published than withdrawn. This paper and its figures are licensed CC BY 4.0
and the accompanying code Apache-2.0.

\bibliographystyle{plainnat}
\bibliography{references}

\appendix

\section{Withdrawals, Corrections and Registered Deviations}
\label{sec:corrections}

This work's earlier drafts published claims it now withdraws. Reporting guidance for pre-registered
work asks that deviations be disclosed in
one systematic place rather than left for a reader to reconstruct by comparing documents
\citep{willroth2024bestlaid}, so every such change affecting \emph{this} paper's claims is
collected here and the body text states only the surviving claim. \textbf{Twenty-six} entries follow.
Substantive withdrawals give the claim as
published and what replaces it; arithmetic corrections are grouped by cause, because the useful
information is the mechanism that produced a class of wrong digits rather than each digit, and the
superseded values themselves are recoverable from the repository history released with this paper.
The series' other two papers carry the entries bearing on their own claims.

\begin{enumerate}\itemsep2pt
  \item \textbf{``No measurable advantage for loop-level optimization.''} The original headline,
    across the original ten runs. Withdrawn: the treatment never reached the gene, so those
    runs compared a genome against itself (\S\ref{sec:defect}). What replaces it is one
    configuration in which the treatment demonstrably ran, reported in the companion paper on the conditional lever.
  \item \textbf{``The null is tight rather than merely underpowered.''} It rested on four
    intervals of exactly $[0,0]$. An interval of $[0,0]$ arises when the two arms produce
    byte-identical genomes and the cassette therefore serves identical completions, which is the
    sharpest available evidence that the experiment did \emph{not} run rather than support for
    its conclusion.
  \item \textbf{``The extra evolvable gene over-fits the train split.''} Untestable from the
    release as published, because the train aggregate was computed and discarded; and the
    default one-component-per-proposal selector predicts the \emph{opposite}, since the two-gene
    arm searched the shared gene about half as hard (defect~3). A later split sweep then found no
    split-size trend to overfit along at all.
  \item \textbf{The retraction accounting.} Published as eleven superseded runs whose zero
    $\dd$ values were all identities. It is \emph{fifteen}, and four of them are not zero
    (\S\ref{sec:superseded}); an unchanged verify gene does not make the two arms the same
    program, because they still differ through the gene they share.
  \item \textbf{``$\dd(D-B) = -0.10$ excludes zero,'' and that it rules out under-powering.} A
    percentile-index error dropped the largest resample atom, which at $n=3$ carries
    $1/27 = 3.7\% > \alpha/2$. Corrected, the interval is $[-0.150,+0.000]$ and includes zero
    (Table~\ref{tab:ciexcl}). The error's direction flatters negative effects.
  \item \textbf{What the two controls establish.} An earlier caption inferred from arm~A's and
    arm~C's flatness that nothing else in the pipeline varies, and therefore that arm~D's drop
    was attributable to optimization. Both are deterministic replays of a fixed genome and bound
    nothing about the variance of a \emph{fresh} draw.
  \item \textbf{``This study contains no measurement noise.''} Published as a scope note, and
    load-bearing. False: ten released runs re-score one fixed genome on one fixed held-out set,
    giving $\sigma(\dd) = 0.0291$ on a thirty-task mean (\S\ref{sec:sigma}). Every effect in
    the series is now reported against it.
  \item \textbf{Six $\sigma$ quotients restated in prose.} $\sigma$ was re-estimated from five
    replicates to ten during the study. One machine-checked sentence per quantity was updated and
    the restatements elsewhere were not, so six quotients continued to divide by the superseded
    value; one of them sat in the same sentence as a number that had been corrected. Two more
    divided a $\dd$ that the paper had already rounded for display, and one sentence asserted
    that all three within-run $\sigma$ values exceed the cross-run estimate when the saturated
    level's does not. All now derive from the current estimate, and a coverage test fails if any
    paper in the series prints a $\sigma$ quotient no artifact produces.
  \item \textbf{The five-replicate $t$-tests used a population standard deviation.} Every
    $K{=}5$ standard error divided the \emph{population} sd of the five draws by $\sqrt{5}$, and
    the result was then compared against critical values from $t(4)$, a distribution that
    exists \emph{because} the denominator is $n-1$. The unbiased estimator inflates $\sigma$ by
    $\sqrt{5/4} = 1.118$, so every $t$ in earlier drafts exceeded its corrected value by
    $11.8\%$ (equivalently, correcting one lowers it by $10.6\%$; we state both because the
    bare phrase ``too large by'' does not say which value is the denominator). The same
    mistake in the cross-run estimate, over $30$ paired tasks, \emph{understated} $\sigma(\dd)$
    as $0.0286$ where the unbiased estimator gives $0.0291$, and scaled every ratio built on it
    by $1.7\%$. Both are corrected throughout, and $0.0291$ is the value divided by here. The
    error was invisible to \texttt{verify\_paper\_claims.py}, which checks the paper against the
    same function that computed it, so paper and checker were wrong together, the same
    shared-error blind spot as the borrowed $\sigma$, one level down.
    It is this paper's own thesis arriving one level down, and worth stating as such: the
    diagnostic we propose catches a treatment that never applied, and would not have caught this.
  \item \textbf{``The ratios reported there are the right ones and stand.''} Said of the
    companion ceiling paper's borrowed-$\sigma$ ratios. Half right: the cross-run estimate is the
    correct denominator \emph{of the two available} for a between-run curve, which is what this
    section establishes. But that paper has since measured each row's own $\sigma$ and superseded
    those ratios, so ``stand'' overstated a choice between two borrowed denominators as a final
    answer.
  \item \textbf{``In the other fifteen \emph{both} arms accepted nothing on every seed.''} Said of
    the runs the acceptance flag catches, and false in eleven of them. Those files predate the
    \texttt{genes\_changed} field, so the census reads an \emph{absent} key as a zero: this
    paper's own failure mode, occurring inside the diagnostic it proposes as the cure. Their released
    scores contradict the claim directly: an arm that accepts nothing returns the seed genome and
    reproduces its own seed score, and eleven of the fifteen have a seed where it does not. Also
    withdrawn from the same sentence: that the flag ``re-derives the original retraction from the
    scores alone'', which reads no score at all, and that the original retraction was about
    acceptance: it was about the \emph{critique} gene, unchanged in $15/15$ winning genomes while
    \texttt{generate} changed in $5/15$. The counts in the sentence were right throughout; the
    mechanism and the inference were not.
  \item \textbf{``Every experiment was pre-registered \ldots\ committed to version control
    before spending.''} True of most, not all, and the shared harness section now says so with the
    exceptions named. Git shows three kinds. The first single-constraint run was produced at
    05:26 UTC on 2026-07-23 and the earliest registration in the repository was committed at
    08:01 the same day, so it predates the practice. The code-domain runs and the actuator
    experiment each appear first in the commit carrying their own results, with no registration
    before them anywhere. And the feedback-actionability pair and the second-worker replication
    line carry registrations committed in the same commit as their results, seventeen to
    twenty-five minutes and two and a half to three and a quarter hours after the runs' own
    \texttt{meta.generated} timestamps. Those registrations may well have been written first, but
    version control cannot show it, and ``committed to version control before spending'' is a
    claim about version control. Table~\ref{tab:results} separately called all ten of the
    original runs pre-registered, which is wrong for run~1 and for runs~9 and~10; it now calls
    them the original ten. Found on the third pre-submission read, by comparing every run's own
    \texttt{meta.generated} against the commit dates of the registrations that name it.
  \item \textbf{``Our first fifteen contrast runs measured the effect of evolving the
    verifier.''} Correct for thirteen of them. The other two are the actuator runs, where
    arm~\textbf{D} evolved \texttt{generate} and an \texttt{actuator} gene with the verifier
    held fixed by design, so their critique-gene identity is expected rather than diagnostic, and
    the selector confound (defect~3) rather than defect~1 is what invalidates their contrast. The
    $15/15$ count is unchanged; what was wrong is the description of what all fifteen set out to
    measure. Found on the final pre-submission read.
  \item \textbf{``Ten of the fifteen reported a \dd\ of exactly zero.''} Correct as a count only
    under a summation order this paper never specified, and it was the wrong predicate regardless.
    \S\ref{sec:defect} defines $\dd = \mathrm{mean}(D) - \mathrm{mean}(B)$; under that definition
    \emph{eleven} of the fifteen are bit-equal to zero. Under the mean of per-seed differences,
    the form the released analysis code computes, and identical in real arithmetic, ten are, and
    the eleventh is $-2.22 \times 10^{-17}$. One run, \texttt{multihard}, sits between them because
    its per-seed differences are $[0, +0.055, 0, 0, -0.055]$ and cancel. So the published ``ten''
    matched the code while the published definition gave eleven, and ``exactly zero'' turned on
    floating-point associativity.
    The surrounding argument needs an \emph{identity} rather than a zero, and that is invariant:
    ten runs produced held-out scores identical on every seed. The body now claims that, and
    \texttt{verify\_paper\_claims.py} pins both counts so a future edit cannot restore a number that
    moves with summation order. This defect is in the class the series says registry-matching
    structurally cannot reach.
  \item \textbf{``Seven more instances of the same class.''} Six, at the time of this correction.
    Table~\ref{tab:defects} then listed nine defects and its caption assigned $1$--$3$ to the
    original headline, the unchanged critique gene and the two wiring choices that caused it,
    so the instances the audit had added were $4$--$9$.
    \S\ref{sec:defect} always said six; the abstract, contribution~3 and the conclusion said seven,
    which would require a tenth row no table supports. The likely arithmetic is $9-2$: the abstract
    and the conclusion both introduce the original as \emph{two} ordinary wiring choices, which is
    true of its causes and leaves out the symptom the caption counts as defect~1. The count is now
    derived from the table's own row count and the range its caption names, at all four sites,
    which is why the table's later growth to eleven rows moved every site with it instead of
    stranding four numbers at nine. That is the difference between a count that is checked and a
    count that was once correct. The count is also restated in an abstract, the one place a
    reader checks first.
  \item \textbf{``Four runs overrode the \$5 ceiling.''} Five did, at four distinct values. The
    shared harness section this paper \texttt{\textbackslash input}s has said
    ``five runs raised it explicitly, at four distinct values'' throughout, and
    \texttt{verify\_paper\_claims.py} checked \emph{that} sentence and agreed with it. This paper
    restated the same fact in its own words two sections later, and the checker's scraper uses a
    first-match search, so the paraphrase was never read. The mechanism is worth more than the
    digit: a fact stated in two wordings is guarded at one site and free to rot at every other, and
    the same slip had already been made once in this pair of numbers, which is why the shared
    sentence is checked at all. Every site stating this count is now checked, in the papers and in
    the repository's documentation, by finding the sites rather than listing them; a sixth site in
    \texttt{docs/EXPERIMENTS.md} was carrying the same wrong digit and was found by that sweep.
    The paragraph reporting it is the one headed
    ``two factual slips belong with the withdrawals'', so the correction carried a third.
  \item \textbf{``The discarded first four-arm attempt at $-0.094$, ${\approx}3.24\sigma$.''}
    Two errors in one clause, and the run was the wrong one. This paper uses
    \emph{the discarded first attempt at the four-arm re-run} twice elsewhere for \texttt{c2}, whose
    mean $\dd$ is $-0.015$ ($0.52\sigma$) and which the pre-correction partition excludes outright,
    since that partition is defined by \texttt{module\_selector} and \texttt{c2} carries the
    corrected value. The run that measures $-0.094$ is run~1, the single-constraint one, which
    Table~\ref{tab:results} lists as \emph{retained}. The count the sentence supports is
    unaffected, since two of the four non-zeros still sit above $1\sigma$, so this is a naming error, not a
    result. The multiple was also wrong: $\sigma$ multiples in this series are the exact $\dd$
    divided by the unrounded $\sigma$, which gives $3.23$, and $3.24$ is what the paper's own
    rounded $-0.094$ over a rounded $0.029$ gives. The check that should have caught it was fed the
    same rounded $0.094$, so both sides of the comparison descended from the display and it passed
    inside tolerance while the artifact's value failed. It now reads the run.
  \item \textbf{An enumeration of the twenty-five replayable runs that sums to twenty-seven.}
    \S\ref{sec:fix} lists them as the two hand-written-gene conditions, the four-arm re-run, the
    two new $p_0$ levels, the two calibration passes, the bought rows of the split and budget
    sweeps, the cross-provider critic run, ``and the eleven runs of the second-worker line''. The
    first seven groups are $16$; the registry's remainder is $9$, not $11$, eight runs on
    \texttt{llama-3.3-70b-versatile} plus the diagonal splice control, which runs a third model
    (\texttt{llama-3.1-8b-instant}) and is therefore not part of that line. The two files that
    make the literal count eleven are a spend report, which is not a run, and the cross-provider
    run, which the same sentence already counts. ``Eleven'' also opens the following sentence, for
    a different quantity, which is where it was most likely copied from, the same
    collide-with-a-neighbour signature as the depth bound in Paper~I's register. \emph{The total
    was right and the parts were not}, which is why no check caught it: twenty-five is pinned
    against the registry, and the enumeration behind it was prose. The parts are now summed
    against \texttt{verify\_reproducibility}'s own replayable count, so the list and the total
    fail together or not at all.
  \item \textbf{``Auditing the fix, and then auditing the audit, found the same class of error
    eight more times.''} Eight is right and the attribution was not: two of the eight arrived
    while replicating on a second worker, which is what Table~\ref{tab:defects}'s caption says
    and what three prose sites contradicted. Every count involved is pinned and all three were
    correct, so the guard could not see it; the claim that moved was about \emph{how} the
    defects were found, and no artifact records that.
  \item \textbf{``The current state-of-the-art reflective optimizer.''} An unscoped superlative
    about an open set at an unstated time, and nothing derived it. This paper's own rule, written
    for the other one, is that a superlative is a claim about a set and has to say which set. The
    sentence now carries GEPA's ICLR-2026 oral instead, which is the credential the superlative
    was standing in for, is checkable in the bibliography, and cannot go stale between submission
    and a rolling review. Nothing quantitative moves.
  \item \textbf{``Every $\dd$ in Table~\ref{tab:results} is a comparison of a genome against
    itself.''} This register already withdrew the identity reading, and three places in the body
    went on asserting it, one of them the opening of the section whose own text calls it too
    strong sixty-four lines later. The table it quantifies over prints two non-zero rows. The
    surviving claim is the one the withdrawal states: the gene that defines the contrast never
    changed, so nothing in that column measures the treatment, and where it is non-zero the arms
    differ in the gene they share. A correction recorded and not carried into the body is this
    paper's own failure mode, one level up.
  \item \textbf{``\dots and never touch the component.''} The released summary records that GEPA
    proposed critique rewrites in the flagship run and that they were rejected, having been
    generation prose written into a critique slot. So the optimizer did reach the component; what
    it never did was change it. The diagnosis is unaffected and the phrasing was false.
  \item \textbf{``The per-gene accepted-change count.''} The harness records which genes of the
    \emph{winning} genome differ from the seed. That is an identity test, not a tally of accepted
    proposals, and an empty count cannot separate a search that proposed and lost from one that
    never proposed. The field, the reporting rule built on it and the strata named after it are
    renamed accordingly.
  \item \textbf{``Restricted to the seeds where arm~D accepted something, $\dd = +0.0100$.''}
    That restriction conditions on an outcome of the treatment rather than on the assignment, so
    the figure is not a treatment-effect estimate. The stratification's value is that it locates
    where the aggregate comes from, and it is now reported as that.
  \item \textbf{``It would have caught defects~1--3 before any money was spent.''} The positive
    control runs to this day with \texttt{module\_selector} at GEPA's default and no verify-truth
    signal, and returns $+1.000$ anyway, so it does not fail under defects~2 or~3; its reflection
    function also returns the same proposal whatever records it is shown. The claim is narrowed to
    defect~1, and the control's use of one task set as both train and held-out is now disclosed.
  \item \textbf{``The $5.9\times$ spread is measured across nine genomes.''} The nine cells are
    arm-by-level aggregates, each pooling three seeds' genomes, across levels that differ in which
    constraints are excluded. The spread stands; what it is measured across does not hold genome,
    seed and task regime fixed, and the sentence said it did.
\end{enumerate}

\end{document}
